\documentclass[12pt,a4paper]{article}

% --- Packages ---
\usepackage[T1]{fontenc}
\usepackage[utf8]{inputenc}
\usepackage[magyar]{babel}
\usepackage{geometry}
\usepackage{booktabs}
\usepackage{array}
\usepackage{longtable}
\usepackage{hyperref}
\usepackage{tabularx}
\usepackage{titlesec}
\usepackage{parskip}
\usepackage{microtype}
\usepackage{lmodern}
\usepackage{xcolor}
\usepackage{fancyhdr}
\usepackage{graphicx}

% --- Page geometry ---
\geometry{
  top=2.5cm,
  bottom=2.5cm,
  left=3cm,
  right=2.5cm
}

% --- Colors ---
\definecolor{accentblue}{RGB}{30, 80, 160}
\definecolor{lightgray}{RGB}{245, 245, 245}
\definecolor{tablegray}{RGB}{230, 235, 245}

% --- Hyperref setup ---
\hypersetup{
  colorlinks=true,
  linkcolor=accentblue,
  urlcolor=accentblue,
  citecolor=accentblue,
  pdftitle={Ipari Képfeldolgozás --- 1. Mérföldkő},
  pdfauthor={Ilonczai András}
}

% --- Section formatting ---
\titleformat{\section}
  {\large\bfseries\color{accentblue}}
  {\thesection.}{0.5em}{}[\titlerule]

\titleformat{\subsection}
  {\normalsize\bfseries\color{accentblue!80!black}}
  {\thesubsection}{0.5em}{}

% --- Header / Footer ---
\pagestyle{fancy}
\fancyhf{}
\rhead{\small Ipari Képfeldolgozás Lab.\ Gyak.\ | 1.\ Mérföldkő}
\lhead{\small Ilonczai András}
\rfoot{\small \thepage}
\renewcommand{\headrulewidth}{0.4pt}

% --- Table column types ---
\newcolumntype{L}[1]{>{\raggedright\arraybackslash}p{#1}}
\newcolumntype{C}[1]{>{\centering\arraybackslash}p{#1}}
\newcolumntype{R}[1]{>{\raggedleft\arraybackslash}p{#1}}

% ================================================================
\begin{document}

% --- Title page block ---
\begin{center}
  {\LARGE\bfseries\color{accentblue} Ipari Képfeldolgozás Laboratóriumi Gyakorlat}\\[0.4em]
  {\large 1.\ Mérföldkő --- Algoritmus Terv \& Költségbecslés}\\[1.2em]
  {\large\bfseries Projekt: Objektumok regisztrációja 3D pontfelhőn}\\[1em]
  \hrule
  \vspace{0.8em}
  \begin{tabular}{ll}
    \textbf{Beadási határidő:} & 2026. március 1.\ 23:59 \\
    \textbf{Csapat tagjai:}    & Ilonczai András \\
    \textbf{Dátum:}            & 2026. február \\
  \end{tabular}
  \vspace{0.8em}
  \hrule
\end{center}

\vspace{1em}

% ================================================================
\section{Projekt áttekintés}

A projekt célja 3D pontfelhők alapján történő objektum-regisztráció és illesztés megvalósítása. A regisztráció során két vagy több, különböző nézőpontból felvett pontfelhőt illesztünk össze egyetlen koherens háromdimenziós modellé. A folyamat alapvetően két lépésből áll: durva (\textit{coarse}) és finom (\textit{fine/dense}) illesztésből.

A pontfelhő-regisztráció számos iparági alkalmazásban kulcsfontosságú, beleértve a minőségellenőrzést, a reverz tervezést (\textit{reverse engineering}) és a robotikai navigációt. A projekt eredménye egy olyan pipeline lesz, amely képes automatikusan illeszteni többszörös szkennelésből származó pontfelhőket.

% ================================================================
\section{Algoritmus terv}

\subsection{Feldolgozási pipeline áttekintése}

A teljes feldolgozási folyamat az alábbi főbb lépésekből áll:

\begin{enumerate}
  \item \textbf{Adatbevitel \& előfeldolgozás} \\
    A nyers pontfelhők beolvasása, zajszűrés (\textit{Statistical Outlier Removal}, \textit{Radius Outlier Removal}), voxel grid alapú ritkítás (\textit{downsampling}) a számítási igény csökkentéséhez.

  \item \textbf{Jellemzők kinyerése (\textit{Feature Extraction})} \\
    Helyi geometriai jellemzők számítása: Normal becslés (PCA alapú), FPFH (\textit{Fast Point Feature Histograms}) vagy SHOT deszkriptorok alkalmazása kulcspontokhoz.

  \item \textbf{Durva illesztés (\textit{Coarse Registration})} \\
    RANSAC alapú kezdeti transzformáció becslése a leíró jellemzők alapján. Célja: megközelítőleg helyes elhelyezés biztosítása a finom illesztés számára.

  \item \textbf{Finom illesztés (\textit{Fine Registration} --- ICP)} \\
    \textit{Iterative Closest Point} (ICP) algoritmus alkalmazása a durva illesztés eredményére. Variánsok: \textit{Point-to-Point} ICP, \textit{Point-to-Plane} ICP. Konvergencia-kritérium: RMSE küszöbérték vagy max.\ iterációs szám.

  \item \textbf{Minőségértékelés} \\
    Az illesztés pontosságának mérése: RMSE (\textit{Root Mean Square Error}), overlap arány, vizuális ellenőrzés. Iteratív finomítás szükség esetén.

  \item \textbf{Eredmény exportálás} \\
    Összefűzött pontfelhő mentése (\texttt{.ply}, \texttt{.pcd} formátumban), transzformációs mátrixok naplózása.
\end{enumerate}

% ----------------------------------------------------------------
\subsection{ICP algoritmus részletezése}

Az ICP (\textit{Iterative Closest Point}) a pontfelhő-regisztráció egyik legelterjedtebb algoritmusa \cite{besl1992}. Az algoritmus iteratívan minimalizálja a forrás- és célpontfelhő közötti távolságot:

\begin{enumerate}
  \item Minden forrásponthoz megkeressük a legközelebbi célpontot (\textit{nearest neighbor search}, pl.\ KD-fa).
  \item Kiszámítjuk az optimális forgatási ($\mathbf{R}$) és eltolási ($\mathbf{t}$) transzformációt (SVD alapú zárt megoldás).
  \item Alkalmazzuk a transzformációt a forráspontfelhőre.
  \item Az RMSE konvergál, vagy elérjük a maximális iterációszámot $\rightarrow$ leállás.
\end{enumerate}

% ----------------------------------------------------------------
\subsection{RANSAC alapú globális regisztráció}

A RANSAC (\textit{Random Sample Consensus}) algoritmust FPFH jellemzőkkel kombinálva alkalmazzuk a kezdeti transzformáció becslésére. Az FPFH jellemzők 33 dimenziós hisztogramban foglalják össze egy pont $k$-szomszédsági geometriáját, ami robusztus egyeztetést tesz lehetővé. A módszert Zhou et al.\ \cite{zhou2016} nyílt forráskódú implementációként is elérhetővé tette az Open3D könyvtárban.

% ----------------------------------------------------------------
\subsection{Felhasznált szoftvereszközök és könyvtárak}

\begin{table}[h!]
\centering
\begin{tabular}{L{3.5cm} C{2cm} L{6.5cm}}
\toprule
\textbf{Eszköz / Könyvtár} & \textbf{Verzió} & \textbf{Felhasználás} \\
\midrule
Python        & 3.10+  & Fő programozási nyelv \\
Open3D        & 0.18   & Pontfelhő I/O, ICP, RANSAC, vizualizáció \\
NumPy         & 1.26   & Numerikus számítások, mátrixműveletek \\
SciPy         & 1.12   & KD-fa, lineáris algebra \\
Matplotlib    & 3.8    & Eredmények vizualizációja, ábrázolás \\
PyVista       & 0.43   & Interaktív 3D vizualizáció \\
CloudCompare  & 2.13   & Referencia-szoftver, manuális ellenőrzés \\
\bottomrule
\end{tabular}
\caption{Felhasznált szoftvereszközök és könyvtárak}
\end{table}

% ================================================================
\section{Szakirodalmi áttekintés}

\begin{thebibliography}{9}
  \bibitem{besl1992}
    Besl, P.\,J., \& McKay, N.\,D. (1992).
    A method for registration of 3-D shapes.
    \textit{IEEE Transactions on Pattern Analysis and Machine Intelligence}, 14(2), 239--256.

  \bibitem{rusu2009}
    Rusu, R.\,B., Blodow, N., \& Beetz, M. (2009).
    Fast point feature histograms (FPFH) for 3D registration.
    \textit{IEEE ICRA}, 3212--3217.

  \bibitem{zhou2016}
    Zhou, Q.-Y., Park, J., \& Koltun, V. (2016).
    Fast global registration.
    \textit{ECCV}, 766--782.

  \bibitem{chen1992}
    Chen, Y., \& Medioni, G. (1992).
    Object modelling by registration of multiple range images.
    \textit{Image and Vision Computing}, 10(3), 145--155.

  \bibitem{zhou2018}
    Zhou, Q.-Y., Park, J., \& Koltun, V. (2018).
    Open3D: A modern library for 3D data processing.
    \textit{arXiv:1801.09847}.
\end{thebibliography}

% ================================================================
\section{Csapattagok és feladatkiosztás}

\begin{table}[h!]
\centering
\begin{tabular}{L{3.5cm} L{7cm} C{2.5cm}}
\toprule
\textbf{Csapattag} & \textbf{Feladat} & \textbf{Mérföldkő} \\
\midrule
Ilonczai András &
  Algoritmus terv, szakirodalom feldolgozás, ICP implementáció.
  Költségbecslés, RANSAC implementáció, eredmény kiértékelés &
  1., 2., 3. \\
\bottomrule
\end{tabular}
\caption{Csapattagok és feladatkiosztás}
\end{table}

% ================================================================
\section{Projekt weboldal}

A projekt haladását egy dedikált weboldalon tartom nyilván, amelyen mérföldkőnként bemutatjuk az elvégzett fejlesztéseket, eredményeket.

\textbf{Weboldal URL:} \texttt{[ide kell a weboldal linkje]}

Az oldalon elérhető tartalmak:
\begin{itemize}
  \item Projekt leírás és célkitűzések
  \item Algoritmus terv és szakirodalom összefoglalása (1.\ mérföldkő)
  \item Kezdeti implementáció és eredmények bemutatása (2.\ mérföldkő)
  \item Végleges eredmények és értékelés (3.\ mérföldkő)
\end{itemize}

% ================================================================
\section{Költségbecslés}

Az alábbiakban összefoglaltam a projekt megvalósításához szükséges hardver- és szoftvereszközök becsült költségeit. Megjegyzés: a laborban rendelkezésre álló eszközöket nem szükséges megvásárolni, az árak tájékoztató jellegűek.

\subsection{Hardverköltségek}

\begin{table}[h!]
\centering
\begin{tabular}{L{4cm} L{4.5cm} C{2.5cm} L{3cm}}
\toprule
\textbf{Eszköz} & \textbf{Leírás} & \textbf{Becsült ár (USD)} & \textbf{Megjegyzés} \\
\midrule
NextEngine 3D Scanner HD &
  Nagy felbontású asztali 3D szkenner, 50--200 mikron pontosság, 2 lézercsík &
  2\,995 &
  Laborban rendelkezésre áll \\
Számítógép (munkaállomás) &
  Intel Core i7/i9, 32 GB RAM, NVIDIA RTX 3060+ GPU &
  1\,500--2\,500 &
  Laborban elérhető \\
Forgóasztal / tartójig &
  Objektum rögzítéséhez szkennelés közben &
  50--150 &
  Laborban elérhető \\
Kalibrációs lap &
  NextEngine kalibráló célpont &
  $\sim$50 &
  Szkennerhez tartozik \\
\midrule
\multicolumn{2}{l}{\textbf{HW részösszeg (ha megvásárolnánk):}} &
  \multicolumn{2}{l}{\textbf{$\sim$4\,595--5\,695 USD}} \\
\bottomrule
\end{tabular}
\caption{Hardverköltségek}
\end{table}

\subsection{Szoftverköltségek}

\begin{table}[h!]
\centering
\begin{tabular}{L{4.5cm} L{3.5cm} C{2.5cm} L{3cm}}
\toprule
\textbf{Szoftver} & \textbf{Licensz típus} & \textbf{Ár / év (USD)} & \textbf{Megjegyzés} \\
\midrule
Python 3.10+              & Nyílt forráskód (PSF)              & 0 & Ingyenes \\
Open3D 0.18               & MIT License                        & 0 & Ingyenes \\
NumPy / SciPy / Matplotlib & BSD License                       & 0 & Ingyenes \\
PyVista                   & MIT License                        & 0 & Ingyenes \\
CloudCompare 2.13         & GNU GPL v2                         & 0 & Ingyenes \\
ScanStudio HD (NextEngine) & Kereskedelmi (szkennerhez jár)    & 0 & Szkenner szoftver \\
Jupyter Notebook / VS Code & MIT / Ingyenes                   & 0 & Fejlesztői környezet \\
\midrule
\multicolumn{2}{l}{\textbf{SW részösszeg:}} &
  \multicolumn{2}{l}{\textbf{0 USD} (minden szoftver ingyenes)} \\
\bottomrule
\end{tabular}
\caption{Szoftverköltségek}
\end{table}

\subsection{Összesített költségbecslés}

\begin{table}[h!]
\centering
\begin{tabular}{L{5.5cm} C{3.5cm} L{4.5cm}}
\toprule
\textbf{Kategória} & \textbf{Becsült költség (USD)} & \textbf{Megjegyzés} \\
\midrule
Hardver (szkenner + PC)            & 4\,595--5\,695 & Laborban rendelkezésre áll \\
Szoftver (összes eszköz)           & 0              & Nyílt forráskódú / ingyenes \\
Egyéb (kalibrációs testek, stb.)   & $\sim$50--100  & Opcionális \\
\midrule
\textbf{ÖSSZESEN (ha nulláról vennénk)} & \textbf{$\sim$4\,645--5\,795 USD} & \\
\textbf{Tényleges projektköltség}       & \textbf{$\sim$0 USD} & Labor infrastruktúra \\
\bottomrule
\end{tabular}
\caption{Összesített költségbecslés}
\end{table}

\vspace{0.5em}
\noindent\textit{Megjegyzés:} A projekt kizárólag nyílt forráskódú szoftvereket és a laboratóriumi infrastruktúrát vesz igénybe, ezért a tényleges projektköltség minimális. Kereskedelmi 3D szkenner szoftver (pl.\ Artec Studio, GeoMagic) alternatívaként szóba jöhetne, de ezek éves licenszdíja 1\,500--5\,000 USD körül van --- ezeket a projekt nem igényli.

% ================================================================
\vspace{2em}
\noindent\rule{\textwidth}{0.4pt}\\
\small\textit{Ipari Képfeldolgozás Lab.\ Gyak.\ | 1.\ Mérföldkő | 2026. február 25.}

\end{document}